\documentclass[12pt,a4paper]{article}

% ------------------------------------------------
% Packages
% ------------------------------------------------
\usepackage[margin=1in]{geometry}
\usepackage{setspace}
\usepackage{enumitem}
\usepackage{booktabs}
\usepackage{titlesec}
\usepackage{hyperref}
\usepackage{fancyhdr}

\setstretch{1.2}
\setlength{\parindent}{0pt}
\setlength{\parskip}{6pt}

% Header & Footer
\pagestyle{fancy}
\fancyhf{}
\fancyhead[L]{Psychology – Paper 3 (HL)}
\fancyhead[R]{Resource Booklet}
\fancyfoot[C]{\thepage}

% Section formatting
\titleformat{\section}{\large\bfseries}{\thesection}{1em}{}
\titleformat{\subsection}{\normalsize\bfseries}{}{0em}{}

% ------------------------------------------------
% Document
% ------------------------------------------------
\begin{document}

\begin{center}
    {\LARGE \textbf{Paper 3 – Resource Booklet (HL)}}\\
    \vspace{0.5em}
    \textbf{Subject:} Psychology\\
    \textbf{Assessment Focus:} Data Analysis and Interpretation\\
    \textbf{Linked Examination:} Paper 3 (HL)\\
    \textbf{Examination Duration:} 1 hour 45 minutes
\end{center}

\vspace{1em}

\hrule
\vspace{1em}

% ------------------------------------------------
\section*{Instructions to Candidates}

\begin{itemize}
    \item This booklet contains the research sources required to answer all questions in Paper 3.
    \item Read all sources carefully before beginning your responses.
    \item All answers must refer directly to the relevant source(s).
    \item Use appropriate psychological terminology throughout.
    \item Where required, integrate your HL extension knowledge with the information provided.
    \item Do not write your answers in this booklet unless instructed.
\end{itemize}

\vspace{1em}
\hrule
\vspace{1em}

% ------------------------------------------------
\section*{Overview of the Sources}

The sources in this booklet relate to \textbf{one HL extension inquiry topic}.  

Possible topics include:

\begin{itemize}
    \item The role of culture in shaping behavior
    \item The role of motivation in shaping behavior
    \item The role of technology in shaping behavior
\end{itemize}

The specific topic assessed is reflected through the research materials provided.

\vspace{1em}

Sources may include both:

\begin{itemize}
    \item Quantitative research findings
    \item Qualitative research findings
\end{itemize}

Research may be experimental or non-experimental.  
Data may be authentic or created for assessment purposes.

\vspace{1em}
\hrule
\vspace{1em}

% ------------------------------------------------
\section*{Structure of the Sources}

\subsection*{Source Labelling}

\begin{itemize}
    \item Each source is clearly labelled (e.g., Source 1, Source 2, Source 3).
    \item Sources vary in length and format.
    \item Some sources may contain graphs, tables, or statistical summaries.
    \item Others may contain interview excerpts, thematic summaries, or methodological descriptions.
\end{itemize}

\subsection*{Quantitative Sources May Include}

\begin{itemize}
    \item Graphs and charts
    \item Tables of numerical data
    \item Statistical outcomes
    \item Experimental or correlational findings
\end{itemize}

\subsection*{Qualitative Sources May Include}

\begin{itemize}
    \item Interview extracts
    \item Observational summaries
    \item Thematic analyses
    \item Researcher commentary
\end{itemize}

\vspace{1em}
\hrule
\vspace{1em}

% ------------------------------------------------
\section*{Use of the Sources in the Examination}

The sources support the four questions in Paper 3:

\begin{itemize}
    \item Interpretation of graphical or numerical data
    \item Analysis of research findings
    \item Evaluation of research considerations (e.g., credibility, bias, transferability)
    \item Synthesis of information across multiple sources
\end{itemize}

Students are expected to:

\begin{itemize}
    \item Refer explicitly to source material
    \item Interpret findings accurately
    \item Identify methodological strengths and limitations
    \item Integrate HL extension knowledge where required
    \item Develop reasoned, evidence-based conclusions
\end{itemize}

\vspace{1em}
\hrule
\vspace{1em}

% ------------------------------------------------
\section*{Preparation Guidance}

To effectively use this resource booklet, students should:

\begin{itemize}
    \item Develop strong skills in data interpretation
    \item Understand both quantitative and qualitative research methods
    \item Be familiar with research terminology
    \item Practise synthesizing findings across multiple sources
    \item Structure analytical responses clearly and logically
\end{itemize}

\vspace{2em}

\begin{center}
\textit{End of Resource Booklet}
\end{center}

\end{document}